\chapter{Haakjes wegwerken}
\label{chap:haakjes}

In de wiskunde kom je regelmatig formules tegen waarin haakjes staan, 
zoals $3(x+4)$ of $-(2x-5)$. Soms is zo'n formule met haakjes handig, 
maar vaak wil je hem juist herschrijven zonder haakjes. Dat noemen 
we \emph{haakjes wegwerken}.

Om haakjes weg te werken vermenigvuldigt je de \emph{factor} die v\'o\'or de haakjes staat met elke \emph{term} binnen de haakjes. \\ 
Een \emph{term} is een onderdeel van een som: alles wat door plus- en mintekens gescheiden wordt: zo heeft $3x^2 - 5x + 7$ drie termen en $2x+4y$ er twee). De \emph{factor} is alles wat door vermenigvuldigingstekens gescheiden wordt ($5x$ bestaat uit de 
factoren $5$ en $x$).\\

Zo geldt bijvoorbeeld:
\[
3(x+4) = 3\cdot x + 3\cdot 4 = 3x + 12
\]

Let daarbij goed op het teken: een minteken v\'o\'or de haakjes klapt de 
tekens binnen de haakjes om:
\[
-(2x-5) = -2x + 5
\]

Haakjes wegwerken is een basisvaardigheid die je in bijna elk vervolgonderwerp 
terugziet, zoals vergelijkingen oplossen, formules herschrijven en het 
werken met kwadratische functies. In dit hoofdstuk oefen je deze vaardigheid 
stap voor stap, eerst met eenvoudige voorbeelden en later met formules met 
meerdere haakjes en machten.

\section{Haakjes wegwerken}
In wiskundige uitdrukkingen komen vaak ``haakjes'' voor. In deze paragraaf komen de rekenregels aan de orde met betrekking tot het wegwerken van haakjes bij vermenigvuldigingen.

\regel{Regel 1}{a \cdot (b + c) = a \cdot b + a \cdot c}

Voorbeelden:

\begin{align*}
4 \cdot (3x + y) &= 12x + 4y \\
5 \cdot (2x + 3y) &= \\
a \cdot (2x + 3y) &= a \cdot 2x + a \cdot 3y 
\end{align*}

Dit antwoord wordt meestal geschreven als: $2ax + 3ay$

\begin{align*}
2p \cdot (3x - 2y) &= \\
6(3x - 2y + 5z) &= 18x - 12y + 30z \\
3(-x + 4y - 8z) &= \\
x(3x + 5) &= 3x^2 + 5x \\
y(5y - 3) &=
\end{align*}


\regel{Regel 2}{(a+b)\cdot(c+d)=a\cdot c+a\cdot d+b\cdot c+b\cdot d}
Voorbeelden:

\begin{align*}
(x + 4) \cdot (x + 2) &= x^2 + 2x + 4x + 8 = x^2 + 6x + 8 \\
(x + 5) \cdot (2x + 3) &= \\
(x - 2y) \cdot (x + y) &= x^2 + x \cdot y - 2x \cdot y - 2y^2 = x^2 - x \cdot y - 2y^2
\end{align*}

Dit antwoord wordt meestal geschreven als: $x^2 - xy - 2y^2$

\begin{align*}
(2x + y)(3x - 4y) &=
\end{align*}

\section{Merkwaardige producten}
Het uitwerken van haakjes kan in sommige gevallen worden versneld door gebruik te maken van de volgende "merkwaardige producten".

\regel{Merkwaardig product 1}{(a+b)(a-b)=a^2-b^2}

Verantwoording van deze algemene regel:
\[(a+b)(a-b)=a^2+ab-ab-b^2=a^2-b^2\]

Voorbeelden:
\begin{align*}
(x + 4)(x - 4) &= x^2 - 16 \\
(y + 3)(y - 3) &= \\
(2x + 5y)(2x - 5y) &= (2x)^2 - (5y)^2 = 4x^2 - 25y^2 \\
(4a - 7b)(4a + 7b) &=
\end{align*}

\regel{Merkwaardig product 2}{(a+b)^2=a^2+2ab+b^2}
\[(a+b)^2=(a+b)(a+b)=a^2+ab+ab+b^2=a^2+2ab+b^2\]

Voorbeelden:
\begin{align*}
(x + 4)^2 &= x^2 + 2 \cdot 4x + 4^2 = x^2 + 8x + 16 \\
(y + 3)^2 &= \\
(2a + 5b)^2 &= (2a)^2 + 2 \cdot 2a \cdot 5b + (5b)^2 = 4a^2 + 20ab + 25b^2 \\
(3x + 2y)^2 &=
\end{align*}

\regel{Merkwaardig product 3}{(a+b)^2=a^2+2ab+b^2}
Verantwoording van deze algemene regel:
\[(a-b)^2=(a-b)(a-b)=a^2-ab-ab+b^2=a^2-2ab+b^2\]

Voorbeelden:
\begin{align*}
(x - 1)^2 &= x^2 - 2 \cdot x \cdot 1 + 1^2 = x^2 - 2x + 1 \\
(y - 6)^2 &= \\
(5p - 2q)^2 &= 25p^2 - 20pq + 4q^2 \\
(4x - 3y)^2 &=
\end{align*}

\newpage
\section{Huiswerkopdrachten}
Hieronder staan twee soorten opgaven. Bij Opgave 1 werk je haakjes weg door vermenigvuldigen. Bij Opgave 2 gebruik je de bekende merkwaardige producten. 

Tip: Schrijf elke tussenstap op zodat je controle kunt houden op je werk.

\subsection*{Opgave 1: Werk de haakjes weg.}
\renewcommand{\arraystretch}{1.6}
\begin{tabular}{m{15em}@{\hspace{4em}}l}
\text{a.} $6(2x - y)$       & \text{g.} $(a + 4)(a - 2)$    \\
\text{b.} $6(3x - 2y)$      & \text{h.} $k(2 - k)$          \\
\text{c.} $2(x - 2y + 3)$   & \text{i.} $(x - y)(2x + y)$   \\
\text{d.} $2a(4x + 5y)$     & \text{j.} $(4p - 1)(2p - 3)$  \\
\text{e.} $p(4p - 5)$       & \text{k.} $(x + 2y)(2x - 3y)$ \\
\text{f.} $4(-2x - y + 5)$  & \\
\end{tabular}

\subsection*{Opgave 2: Werk de haakjes weg door gebruik te maken van merkwaardige producten.}
\renewcommand{\arraystretch}{1.6}
\begin{tabular}{m{15em}@{\hspace{4em}}l}
\text{a.} $(x + 5)(x - 5)$      & \text{f.} $(4y - 3)^2$  \\
\text{b.} $(x - 6)(x + 6)$      & \text{g.} $(2y + 5)^2$  \\
\text{c.} $(2m + 3)(2m - 3)$    & \text{h.} $(2a + 6p)^2$ \\
\text{d.} $(4x + 3y)(4x - 3y)$  & \text{i.} $(a - 3)^2$   \\
\text{e.} $(p + 5)^2$           & \text{j.} $(4x - 7)^2$  \\
\end{tabular}